\documentclass[12pt, a4paper]{report}

% ── Encodage & langue ────────────────────────────────────────────────────────
\usepackage[utf8]{inputenc}
\usepackage[T1]{fontenc}
\usepackage[french]{babel}

% ── Mise en page ─────────────────────────────────────────────────────────────
\usepackage[top=2.5cm, bottom=2.5cm, left=2.5cm, right=2.5cm]{geometry}
\usepackage{setspace}
\onehalfspacing

% ── Polices & couleurs ───────────────────────────────────────────────────────
\usepackage{lmodern}
\usepackage[dvipsnames]{xcolor}
\definecolor{primary}{RGB}{124, 58, 237}   % violet TravelShare
\definecolor{secondary}{RGB}{91, 33, 182}
\definecolor{lightgray}{RGB}{245, 245, 245}
\definecolor{darkgray}{RGB}{70, 70, 70}
\definecolor{success}{RGB}{34, 197, 94}
\definecolor{warning}{RGB}{251, 146, 60}
\definecolor{danger}{RGB}{239, 68, 68}

% ── Graphiques & figures ─────────────────────────────────────────────────────
\usepackage{graphicx}
\usepackage{float}
\usepackage{caption}
\usepackage{subcaption}
\usepackage{tikz}
\usetikzlibrary{shapes.geometric, arrows.meta, positioning, shadows}

% ── Tableaux ─────────────────────────────────────────────────────────────────
\usepackage{booktabs}
\usepackage{tabularx}
\usepackage{array}
\usepackage{multirow}
\usepackage{colortbl}

% ── Listes & énumérations ────────────────────────────────────────────────────
\usepackage{enumitem}
\usepackage{pifont}
\newcommand{\cmark}{\textcolor{success}{\ding{51}}}
\newcommand{\xmark}{\textcolor{danger}{\ding{55}}}
\newcommand{\pmark}{\textcolor{warning}{\ding{119}}}

% ── Code source ──────────────────────────────────────────────────────────────
\usepackage{listings}
\lstset{
  basicstyle=\ttfamily\footnotesize,
  backgroundcolor=\color{lightgray},
  frame=single,
  rulecolor=\color{gray},
  breaklines=true,
  showstringspaces=false,
  keywordstyle=\color{primary}\bfseries,
  commentstyle=\color{gray}\itshape,
  stringstyle=\color{secondary},
  numbers=left,
  numberstyle=\tiny\color{gray},
  tabsize=4
}

% ── Liens & PDF ──────────────────────────────────────────────────────────────
\usepackage{hyperref}
\hypersetup{
  colorlinks=true,
  linkcolor=primary,
  urlcolor=primary,
  citecolor=primary,
  pdftitle={Rapport de mi-parcours -- TravelShare},
  pdfauthor={Mahmoud}
}

% ── En-têtes et pieds de page ────────────────────────────────────────────────
\usepackage{fancyhdr}
\pagestyle{fancy}
\fancyhf{}
\renewcommand{\headrulewidth}{0.4pt}
\fancyhead[L]{\textcolor{primary}{\textbf{TravelShare}}}
\fancyhead[R]{\textit{Rapport de mi-parcours}}
\fancyfoot[C]{\thepage}

% ── Boîtes colorées ──────────────────────────────────────────────────────────
\usepackage{tcolorbox}
\tcbuselibrary{skins, breakable}

\newtcolorbox{infobox}[1][]{
  colback=blue!5!white, colframe=primary,
  fonttitle=\bfseries, title=#1,
  boxrule=0.8pt, arc=4pt
}

\newtcolorbox{warningbox}[1][]{
  colback=orange!10!white, colframe=warning,
  fonttitle=\bfseries, title=#1,
  boxrule=0.8pt, arc=4pt
}

\newtcolorbox{successbox}[1][]{
  colback=green!5!white, colframe=success,
  fonttitle=\bfseries, title=#1,
  boxrule=0.8pt, arc=4pt
}

% ── Commandes personnalisées ──────────────────────────────────────────────────
\newcommand{\placeholder}[2]{%
  \begin{figure}[H]
    \centering
    \begin{tikzpicture}
      \node[draw=gray, dashed, fill=lightgray, minimum width=#1, minimum height=6cm,
            font=\small\itshape\color{darkgray}, align=center, rounded corners=4pt] {#2};
    \end{tikzpicture}
    \caption{#2}
  \end{figure}
}

\newcommand{\todo}[1]{\textcolor{danger}{\textbf{[TODO: #1]}}}

% ─────────────────────────────────────────────────────────────────────────────
\begin{document}

% ── PAGE DE TITRE ─────────────────────────────────────────────────────────────
\begin{titlepage}
  \centering
  \vspace*{1cm}

  % Logo ou bandeau
  \begin{tikzpicture}
    \node[fill=primary, minimum width=14cm, minimum height=2.5cm, rounded corners=8pt,
          font=\Huge\bfseries\color{white}] {\texttt{travel\!\!}\textbf{$\bullet$}\texttt{share}};
  \end{tikzpicture}

  \vspace{1.5cm}
  {\Huge \textbf{Rapport de mi-parcours}}\\[0.4cm]
  {\Large \textcolor{primary}{Projet Android -- \textit{Traveling}}}\\[0.3cm]
  {\large Partie 1 : TravelShare}

  \vspace{1cm}
  \rule{10cm}{0.5pt}\\[1cm]

  \begin{tabular}{rl}
    \textbf{Auteur}       & Mahmoud \\[0.3cm]
    \textbf{Date}         & 27 avril 2026 \\[0.3cm]
    \textbf{Module}       & Développement Android \\[0.3cm]
    \textbf{Dépôt}        & \texttt{travelshare-android} \\
  \end{tabular}

  \vfill
  \textcolor{gray}{\small Application Android de partage et découverte de photos de voyages}
\end{titlepage}

% ── TABLE DES MATIÈRES ────────────────────────────────────────────────────────
\tableofcontents
\newpage

% ─────────────────────────────────────────────────────────────────────────────
\chapter{Introduction et contexte du projet}
% ─────────────────────────────────────────────────────────────────────────────

\section{Présentation générale}

Le projet \textbf{Traveling} est une application mobile Android destinée aux voyageurs. Il est
divisé en deux parties principales développées en binôme~:

\begin{itemize}
  \item \textbf{TravelShare (Partie~1)} : application de partage et de découverte de photos
        de voyages -- dont la réalisation m'est confiée.
  \item \textbf{TravelPath (Partie~2)} : générateur intelligent de parcours de visite,
        développé par le second membre du binôme.
  \item \textbf{Passerelle (Partie~3)} : intégration cohérente des deux modules en une
        seule application unifiée (\textit{Traveling}).
\end{itemize}

Ce rapport couvre l'état d'avancement de la \textbf{partie TravelShare} à mi-parcours du
développement.

\section{Objectifs de TravelShare}

TravelShare est une application permettant à des voyageurs de partager, découvrir et interagir
autour de photos de voyage. Elle propose deux modes de fonctionnement~:

\begin{infobox}[Mode anonyme]
Sans compte utilisateur, le visiteur peut~: parcourir les photos publiques, effectuer des
recherches (textuelles ou vocales), consulter les détails d'une photo (lieu, date,
commentaires, itinéraire), aimer ou signaler un contenu.
\end{infobox}

\begin{infobox}[Mode connecté]
Après inscription et connexion, l'utilisateur peut en plus~: publier des photos, les partager
avec des groupes ou les rendre publiques, enrichir les photos (texte, tags, annotations IA),
recevoir des notifications et accéder à un profil personnel.
\end{infobox}

% ─────────────────────────────────────────────────────────────────────────────
\chapter{Architecture technique}
% ─────────────────────────────────────────────────────────────────────────────

\section{Choix technologiques}

L'application est développée en \textbf{Java} natif pour Android, en respectant les bonnes
pratiques recommandées par Google.

\subsection{Patron de conception MVVM}

L'architecture retenue est \textbf{Model-View-ViewModel (MVVM)}, soutenue par les bibliothèques
Jetpack de Google~:

\begin{itemize}
  \item \textbf{ViewModel} : séparation logique métier / couche d'affichage, survie aux
        rotations d'écran.
  \item \textbf{LiveData} : observabilité réactive des données, mise à jour automatique de
        l'interface.
  \item \textbf{Fragment} : navigation modulaire (HomeFragment, SearchFragment, ProfileFragment).
\end{itemize}

\subsection{Bibliothèques utilisées}

\begin{table}[H]
\centering
\caption{Principales dépendances du projet}
\begin{tabularx}{\textwidth}{lXl}
\toprule
\rowcolor{primary!15}
\textbf{Bibliothèque} & \textbf{Rôle} & \textbf{Version} \\
\midrule
Retrofit2 + OkHttp      & Client HTTP REST, communication avec le backend & 2.x \\
Gson Converter          & Désérialisation JSON $\leftrightarrow$ Java POJO & -- \\
Glide                   & Chargement et mise en cache d'images & 4.16.0 \\
Google Maps SDK         & Affichage carte, marqueurs, sélection de lieu & 19.0.0 \\
Firebase Messaging      & Notifications push (FCM) & BOM \\
PhotoView               & Zoom sur image en plein écran & 2.3.0 \\
Cloudinary              & Hébergement et upload d'images dans le cloud & -- \\
Material Design         & Composants UI (BottomSheet, ChipGroup, Badge) & -- \\
FusedLocationProvider   & Géolocalisation de l'appareil & -- \\
\bottomrule
\end{tabularx}
\end{table}

\subsection{Communication réseau}

Le client HTTP est encapsulé dans la classe \texttt{RetrofitClient}, configurée avec~:
\begin{itemize}
  \item Une URL de base dynamique (modifiable depuis l'écran \texttt{IpConfigActivity}
        pour faciliter les tests en réseau local).
  \item Un intercepteur OkHttp qui injecte automatiquement le token JWT (\texttt{Bearer})
        dans chaque requête sortante lorsque l'utilisateur est connecté.
  \item Des délais d'expiration de 60~secondes (connexion, lecture, écriture).
  \item Un intercepteur de journalisation pour le débogage.
\end{itemize}

\subsection{Structure des packages}

\begin{lstlisting}[language=Java, caption=Organisation des packages Java]
com.example.travelshare/
├── api/                    # Ancien client API (remplacé par network/)
├── data/                   # Données mock (tests)
├── model/                  # POJOs (Photo, Group, Comment, Notification)
│   └── response/           # Objets de réponse API (AuthResponse, etc.)
├── network/                # RetrofitClient, ApiService, TokenManager, FCM
├── ui/
│   ├── auth/               # AuthDialog (connexion / inscription)
│   ├── detail/             # PhotoDetailActivity, PhotoDetailViewModel
│   ├── home/               # HomeFragment, adaptateurs, viewmodels, dialogs
│   ├── profile/            # ProfileFragment
│   └── search/             # SearchFragment, SearchViewModel
├── IpConfigActivity.java   # Configuration IP du serveur
├── MainActivity.java       # Point d'entrée, navigation bas de page
├── MainViewModel.java      # État global (connexion, badge notifs)
└── TravelShareApplication.java
\end{lstlisting}

% ─────────────────────────────────────────────────────────────────────────────
\chapter{Fonctionnalités réalisées}
% ─────────────────────────────────────────────────────────────────────────────

\section{Authentification}

\subsection{Connexion et inscription}

Le module d'authentification est implémenté dans \texttt{AuthDialog}, une boîte de dialogue
modale à deux onglets (Connexion / Inscription).

\begin{itemize}
  \item \textbf{Connexion}~: envoi de l'e-mail et du mot de passe vers \texttt{POST /api/auth/login}.
        En cas de succès, le token JWT et les informations utilisateur (nom complet,
        pseudo, identifiant MongoDB) sont persistés localement via \texttt{TokenManager}
        (SharedPreferences chiffré).
  \item \textbf{Inscription}~: collecte du nom complet, du pseudo, de l'e-mail et du mot de
        passe, validation côté client (champs non vides, longueur minimale du mot de passe),
        puis appel à \texttt{POST /api/auth/register}.
  \item L'interface désactive les boutons pendant l'appel réseau pour prévenir les
        soumissions multiples.
\end{itemize}

\begin{figure}[H]
  \centering
  % ── INSÉRER ICI LA CAPTURE DE L'ÉCRAN DE CONNEXION ──
  \begin{tikzpicture}
    \node[draw=gray, dashed, fill=lightgray, minimum width=5cm, minimum height=9cm,
          font=\small\itshape\color{darkgray}, align=center, rounded corners=4pt]
          {Capture d'écran~:\\[0.3cm]\textbf{Dialogue d'authentification}\\(onglets Connexion / Inscription)};
  \end{tikzpicture}
  \caption{Dialogue d'authentification -- onglets Connexion et Inscription}
  \label{fig:auth}
\end{figure}

\subsection{Gestion de session}

La classe \texttt{TokenManager} centralise la persistance de session~:
\begin{itemize}
  \item Sauvegarde du token JWT, du nom complet, du pseudo, de l'identifiant et de la date
        de création du compte.
  \item Vérification de l'état de connexion (\texttt{isLoggedIn()}) utilisée dans
        l'ensemble de l'application pour adapter dynamiquement l'interface.
  \item Invalidation complète à la déconnexion (\texttt{clear()}).
\end{itemize}

L'état de connexion est exposé via un \texttt{LiveData<Boolean>} dans \texttt{MainViewModel},
ce qui permet à toutes les vues de s'adapter automatiquement (affichage du bouton de
connexion, visibilité des FAB, etc.).

% ────────────────────────────────────
\section{Écran d'accueil -- Fil de photos}

\texttt{HomeFragment} est l'écran central de l'application. Il intègre les éléments suivants~:

\subsection{Chargement et affichage des photos}

Les photos publiques sont chargées depuis \texttt{GET /api/photos} (page~1, limite~100) par
\texttt{PhotoViewModel}. La liste est stockée dans un \texttt{MutableLiveData<List<Photo>>}
et filtrée localement à chaque changement de critère, sans nouvel appel réseau.

Trois modes d'affichage sont disponibles, commutables à la volée~:

\begin{itemize}
  \item \textbf{Vue grille}~: \texttt{GridLayoutManager} 2 colonnes, bannière anonyme
        occupant toute la largeur (span~2).
  \item \textbf{Vue liste}~: \texttt{LinearLayoutManager}, cards avec plus d'informations.
  \item \textbf{Vue carte}~: \texttt{GoogleMap} avec marqueurs positionnés sur les
        coordonnées GPS de chaque photo~; un clic sur le marqueur ouvre un
        \texttt{BottomSheet} de prévisualisation.
\end{itemize}

Un \texttt{SwipeRefreshLayout} permet le rechargement manuel du fil.

\begin{figure}[H]
  \centering
  \begin{subcaptionbox}{Vue grille\label{fig:home_grid}}[0.32\textwidth]
    \begin{tikzpicture}
      \node[draw=gray, dashed, fill=lightgray, minimum width=3.8cm, minimum height=8cm,
            font=\tiny\itshape\color{darkgray}, align=center, rounded corners=4pt]
            {Capture d'écran\\vue grille};
    \end{tikzpicture}
  \end{subcaptionbox}
  \hfill
  \begin{subcaptionbox}{Vue liste\label{fig:home_list}}[0.32\textwidth]
    \begin{tikzpicture}
      \node[draw=gray, dashed, fill=lightgray, minimum width=3.8cm, minimum height=8cm,
            font=\tiny\itshape\color{darkgray}, align=center, rounded corners=4pt]
            {Capture d'écran\\vue liste};
    \end{tikzpicture}
  \end{subcaptionbox}
  \hfill
  \begin{subcaptionbox}{Vue carte\label{fig:home_map}}[0.32\textwidth]
    \begin{tikzpicture}
      \node[draw=gray, dashed, fill=lightgray, minimum width=3.8cm, minimum height=8cm,
            font=\tiny\itshape\color{darkgray}, align=center, rounded corners=4pt]
            {Capture d'écran\\vue carte};
    \end{tikzpicture}
  \end{subcaptionbox}
  \caption{Les trois modes de visualisation des photos}
  \label{fig:home_views}
\end{figure}

\subsection{Recherche textuelle et vocale}

\begin{itemize}
  \item \textbf{Recherche par saisie}~: un champ \texttt{EditText} déclenche un filtre
        local en temps réel (via \texttt{TextWatcher}) sur le nom du lieu, le pays, la
        description, le nom de l'auteur et les tags.
  \item \textbf{Recherche vocale}~: le bouton microphone (icône \texttt{ic\_mic}) lance
        l'\texttt{Intent} \texttt{RecognizerIntent.ACTION\_RECOGNIZE\_SPEECH}. Le texte
        reconnu est injecté directement dans le champ de recherche, déclenchant ainsi le
        filtre. La permission \texttt{RECORD\_AUDIO} est demandée au runtime.
\end{itemize}

\subsection{Filtres avancés}

Le bouton \emph{Filtres} ouvre \texttt{FiltersDialog}, une boîte de dialogue proposant~:

\begin{table}[H]
\centering
\caption{Critères de filtrage disponibles}
\begin{tabularx}{\textwidth}{lX}
\toprule
\rowcolor{primary!15}
\textbf{Critère} & \textbf{Implémentation} \\
\midrule
Type de lieu    & Spinner (Nature, Plage, Urbain, Montagne, Musée, \ldots) \\
Auteur          & Spinner alimenté dynamiquement avec les auteurs des photos chargées \\
Période         & Deux DatePickerDialog (date début / date fin), comparaison avec \texttt{createdAt} \\
Rayon géo.      & Latitude, longitude (saisie ou sélection sur carte via \texttt{MapPickerActivity}) + rayon en km~; calcul de distance Haversine \\
Photo similaire & Saisie d'un identifiant de photo (champ prévu, traitement partiel) \\
\bottomrule
\end{tabularx}
\end{table}

\begin{figure}[H]
  \centering
  \begin{tikzpicture}
    \node[draw=gray, dashed, fill=lightgray, minimum width=7cm, minimum height=9cm,
          font=\small\itshape\color{darkgray}, align=center, rounded corners=4pt]
          {Capture d'écran~:\\[0.3cm]\textbf{Dialogue des filtres avancés}};
  \end{tikzpicture}
  \caption{Dialogue des filtres avancés}
  \label{fig:filters}
\end{figure}

% ────────────────────────────────────
\section{Fiche détaillée d'une photo}

\texttt{PhotoDetailActivity} affiche l'ensemble des informations associées à une photo~:

\begin{itemize}
  \item \textbf{Image} en haute résolution (chargée via Glide), cliquable pour accéder à
        la vue plein écran avec zoom (\texttt{FullScreenImageActivity} + PhotoView).
  \item \textbf{Métadonnées}~: auteur (avec avatar-initiale), date, type de lieu, localisation
        (ville + pays).
  \item \textbf{Description} et \textbf{how to get there} (comment s'y rendre).
  \item \textbf{Compteur de likes} et bouton Like/Unlike (bascule en rouge/gris).
  \item \textbf{Bouton Partager} : \texttt{Intent.ACTION\_SEND} vers les applications
        tierces installées.
  \item \textbf{Bouton Signaler} : dialogue de confirmation (signalement logique, pas encore
        d'endpoint API dédié).
  \item \textbf{Bouton Itinéraire} (\emph{Ouvrir dans Maps}) : construction d'un URI
        \texttt{geo:} avec coordonnées et nom du lieu, lancement de Google Maps via Intent ;
        repli sur le navigateur web si Maps n'est pas installé.
  \item \textbf{Section commentaires} : chargement asynchrone, ajout et suppression par
        l'auteur (mode connecté uniquement).
\end{itemize}

\begin{figure}[H]
  \centering
  \begin{subcaptionbox}{Fiche photo\label{fig:detail_main}}[0.48\textwidth]
    \begin{tikzpicture}
      \node[draw=gray, dashed, fill=lightgray, minimum width=5.5cm, minimum height=10cm,
            font=\small\itshape\color{darkgray}, align=center, rounded corners=4pt]
            {Capture d'écran~:\\fiche détaillée de photo};
    \end{tikzpicture}
  \end{subcaptionbox}
  \hfill
  \begin{subcaptionbox}{Plein écran + zoom\label{fig:detail_fullscreen}}[0.48\textwidth]
    \begin{tikzpicture}
      \node[draw=gray, dashed, fill=lightgray, minimum width=5.5cm, minimum height=10cm,
            font=\small\itshape\color{darkgray}, align=center, rounded corners=4pt]
            {Capture d'écran~:\\affichage plein écran avec zoom};
    \end{tikzpicture}
  \end{subcaptionbox}
  \caption{Fiche détaillée d'une photo et vue plein écran}
  \label{fig:photo_detail}
\end{figure}

% ────────────────────────────────────
\section{Publication de photos}

\texttt{PublishBottomSheet} est accessible via le bouton flottant (FAB) visible uniquement
en mode connecté. Le flux de publication se déroule en deux étapes~:

\begin{enumerate}
  \item \textbf{Upload de l'image}~: sélection depuis la galerie Android, copie dans le
        cache de l'application, envoi multipart vers \texttt{POST /api/photos/upload}
        (stockage Cloudinary). L'URL Cloudinary renvoyée est sauvegardée.
  \item \textbf{Création de la fiche photo}~: appel \texttt{POST /api/photos} avec tous
        les métadonnées saisies par l'utilisateur.
\end{enumerate}

\textbf{Champs du formulaire de publication~:}
\begin{itemize}
  \item Description libre.
  \item Lieu (ville) et pays.
  \item Type de lieu (spinner~: Nature, Urbain, Plage, Montagne, Musée, Magasin, Restaurant).
  \item Tags (ajout par chips avec possibilité de suppression individuelle).
  \item Coordonnées GPS (saisie manuelle ou sélection sur carte via \texttt{MapPickerActivity}).
  \item \emph{Comment s'y rendre} (texte libre).
  \item \textbf{Visibilité}~: Public (accessible à tous) ou Privé (réservé à des groupes
        sélectionnés via des cases à cocher).
\end{itemize}

\begin{figure}[H]
  \centering
  \begin{subcaptionbox}{Formulaire de publication\label{fig:publish_form}}[0.48\textwidth]
    \begin{tikzpicture}
      \node[draw=gray, dashed, fill=lightgray, minimum width=5.5cm, minimum height=10cm,
            font=\small\itshape\color{darkgray}, align=center, rounded corners=4pt]
            {Capture d'écran~:\\bottom sheet de publication};
    \end{tikzpicture}
  \end{subcaptionbox}
  \hfill
  \begin{subcaptionbox}{Sélection de lieu sur carte\label{fig:map_picker}}[0.48\textwidth]
    \begin{tikzpicture}
      \node[draw=gray, dashed, fill=lightgray, minimum width=5.5cm, minimum height=10cm,
            font=\small\itshape\color{darkgray}, align=center, rounded corners=4pt]
            {Capture d'écran~:\\MapPickerActivity};
    \end{tikzpicture}
  \end{subcaptionbox}
  \caption{Publication d'une photo et sélection du lieu}
  \label{fig:publish}
\end{figure}

% ────────────────────────────────────
\section{Gestion des groupes}

\texttt{GroupsBottomSheet} propose deux onglets~:

\paragraph{Mes groupes}
Affichage des groupes dont l'utilisateur est membre. Un appui long ouvre un menu contextuel
permettant de \emph{quitter} le groupe ou, si l'utilisateur en est le créateur, de le
\emph{supprimer}. La création d'un nouveau groupe (nom, description, image de couverture)
est accessible depuis cet onglet.

\paragraph{Découvrir}
Liste des groupes existants auxquels l'utilisateur n'appartient pas encore, avec un bouton
\emph{Rejoindre}.

Un clic sur un groupe ouvre \texttt{GroupPhotosActivity}, qui affiche les photos partagées
au sein de ce groupe.

\begin{figure}[H]
  \centering
  \begin{subcaptionbox}{Mes groupes\label{fig:groups_my}}[0.32\textwidth]
    \begin{tikzpicture}
      \node[draw=gray, dashed, fill=lightgray, minimum width=3.8cm, minimum height=8cm,
            font=\tiny\itshape\color{darkgray}, align=center, rounded corners=4pt]
            {Mes groupes};
    \end{tikzpicture}
  \end{subcaptionbox}
  \hfill
  \begin{subcaptionbox}{Découvrir des groupes\label{fig:groups_discover}}[0.32\textwidth]
    \begin{tikzpicture}
      \node[draw=gray, dashed, fill=lightgray, minimum width=3.8cm, minimum height=8cm,
            font=\tiny\itshape\color{darkgray}, align=center, rounded corners=4pt]
            {Découvrir des groupes};
    \end{tikzpicture}
  \end{subcaptionbox}
  \hfill
  \begin{subcaptionbox}{Photos d'un groupe\label{fig:group_photos}}[0.32\textwidth]
    \begin{tikzpicture}
      \node[draw=gray, dashed, fill=lightgray, minimum width=3.8cm, minimum height=8cm,
            font=\tiny\itshape\color{darkgray}, align=center, rounded corners=4pt]
            {Photos du groupe};
    \end{tikzpicture}
  \end{subcaptionbox}
  \caption{Écrans de gestion des groupes}
  \label{fig:groups}
\end{figure}

% ────────────────────────────────────
\section{Notifications}

\subsection{Notifications push (Firebase Cloud Messaging)}

Le service \texttt{MyFirebaseMessagingService} étend \texttt{FirebaseMessagingService}~:
\begin{itemize}
  \item À chaque renouvellement du token FCM, ce dernier est envoyé au backend via
        \texttt{PUT /api/auth/fcm-token} (si l'utilisateur est connecté).
  \item À la réception d'un message push, une notification système est créée avec
        \texttt{NotificationCompat.Builder} (canal haute priorité, redirection vers
        \texttt{MainActivity}).
\end{itemize}

\subsection{Centre de notifications in-app}

\texttt{NotificationsBottomSheet} affiche l'historique des notifications en base~:
\begin{itemize}
  \item Chargement depuis \texttt{GET /api/notifications} et \texttt{GET /api/notifications/unread-count}.
  \item Un badge numérique sur l'icône de la barre de navigation indique le nombre
        de notifications non lues (observable via \texttt{MainViewModel}).
  \item Un clic sur une notification la marque comme lue (\texttt{PATCH /api/notifications/\{id\}/read})
        et redirige vers la ressource concernée (photo ou groupe).
\end{itemize}

\begin{figure}[H]
  \centering
  \begin{tikzpicture}
    \node[draw=gray, dashed, fill=lightgray, minimum width=6cm, minimum height=9cm,
          font=\small\itshape\color{darkgray}, align=center, rounded corners=4pt]
          {Capture d'écran~:\\[0.3cm]\textbf{Centre de notifications}\\(badge + liste)};
  \end{tikzpicture}
  \caption{Centre de notifications in-app avec badge}
  \label{fig:notifications}
\end{figure}

% ────────────────────────────────────
\section{Profil utilisateur}

\texttt{ProfileFragment} s'adapte à l'état de connexion~:
\begin{itemize}
  \item \textbf{Non connecté}~: message d'invitation avec bouton de connexion.
  \item \textbf{Connecté}~: appel à \texttt{GET /api/auth/me} pour obtenir les données
        fraîches depuis le serveur (nom complet, pseudo, date de création du compte).
        Affichage de l'avatar-initiale et de l'année d'inscription.
        Bouton de déconnexion.
\end{itemize}

\begin{figure}[H]
  \centering
  \begin{subcaptionbox}{Non connecté\label{fig:profile_out}}[0.48\textwidth]
    \begin{tikzpicture}
      \node[draw=gray, dashed, fill=lightgray, minimum width=5.5cm, minimum height=9cm,
            font=\small\itshape\color{darkgray}, align=center, rounded corners=4pt]
            {Capture d'écran~:\\profil non connecté};
    \end{tikzpicture}
  \end{subcaptionbox}
  \hfill
  \begin{subcaptionbox}{Connecté\label{fig:profile_in}}[0.48\textwidth]
    \begin{tikzpicture}
      \node[draw=gray, dashed, fill=lightgray, minimum width=5.5cm, minimum height=9cm,
            font=\small\itshape\color{darkgray}, align=center, rounded corners=4pt]
            {Capture d'écran~:\\profil utilisateur connecté};
    \end{tikzpicture}
  \end{subcaptionbox}
  \caption{Fragment de profil utilisateur}
  \label{fig:profile}
\end{figure}

% ────────────────────────────────────
\section{Modèle de données}

Le modèle \texttt{Photo} est le cœur applicatif. Il expose~:

\begin{lstlisting}[language=Java, caption=Extrait du modèle Photo.java]
public class Photo {
    @SerializedName("_id")   private String mongoId;
    private String imageUrl, location, country, locationType;
    private double latitude, longitude;
    private String description, howToGetThere;
    @SerializedName("createdAt") private String date;
    @SerializedName("likes")     private List<String> likesArray;
    private boolean likedByMe, isPublic;
    private List<String> tags;
    private Author author;   // sous-objet { _id, fullName, username, avatar }
}
\end{lstlisting}

% ─────────────────────────────────────────────────────────────────────────────
\chapter{Tableau de bord des fonctionnalités}
% ─────────────────────────────────────────────────────────────────────────────

Le tableau ci-dessous synthétise l'état d'avancement de chaque fonctionnalité spécifiée.

\begin{table}[H]
\centering
\caption{État d'avancement des fonctionnalités TravelShare}
\renewcommand{\arraystretch}{1.4}
\begin{tabularx}{\textwidth}{lXc}
\toprule
\rowcolor{primary!20}
\textbf{Catégorie} & \textbf{Fonctionnalité} & \textbf{État} \\
\midrule
\multirow{4}{*}{\textbf{Authentification}}
  & Inscription (nom, pseudo, e-mail, mot de passe) & \cmark \\
  & Connexion par e-mail / mot de passe              & \cmark \\
  & Persistance de session (JWT + SharedPreferences) & \cmark \\
  & Déconnexion                                      & \cmark \\
\midrule
\multirow{5}{*}{\textbf{Parcours photos}}
  & Chargement du fil public (API REST)              & \cmark \\
  & Vue grille                                       & \cmark \\
  & Vue liste                                        & \cmark \\
  & Vue carte (Google Maps + marqueurs)              & \cmark \\
  & Pull-to-refresh                                  & \cmark \\
\midrule
\multirow{6}{*}{\textbf{Recherche \& filtres}}
  & Recherche textuelle (temps réel)                 & \cmark \\
  & Recherche vocale (RecognizerIntent)              & \cmark \\
  & Filtre par type de lieu                          & \cmark \\
  & Filtre par auteur                                & \cmark \\
  & Filtre par période (date début / fin)            & \cmark \\
  & Filtre par rayon géographique (Haversine)        & \cmark \\
  & Filtre par similarité (photos semblables)        & \pmark \\
\midrule
\multirow{6}{*}{\textbf{Fiche photo}}
  & Affichage image + métadonnées complètes          & \cmark \\
  & Vue plein écran avec zoom (PhotoView)            & \cmark \\
  & Like / Unlike                                    & \cmark \\
  & Commentaires (lecture, ajout, suppression)       & \cmark \\
  & Partage via Intent Android                       & \cmark \\
  & Signalement (dialogue UI)                        & \pmark \\
  & Itinéraire vers le lieu (Intent Google Maps)     & \cmark \\
\midrule
\multirow{5}{*}{\textbf{Publication}}
  & Sélection photo depuis la galerie                & \cmark \\
  & Upload vers Cloudinary (multipart)               & \cmark \\
  & Formulaire complet (lieu, tags, coordonnées…)    & \cmark \\
  & Sélection du lieu sur carte (MapPickerActivity)  & \cmark \\
  & Visibilité public / groupes                      & \cmark \\
  & Suppression de sa propre photo                   & \cmark \\
  & Annotations assistées par IA                     & \xmark \\
\midrule
\multirow{5}{*}{\textbf{Groupes}}
  & Création d'un groupe                             & \cmark \\
  & Rejoindre / quitter un groupe                    & \cmark \\
  & Supprimer un groupe (créateur)                   & \cmark \\
  & Découverte des groupes publics                   & \cmark \\
  & Photos partagées dans un groupe                  & \cmark \\
\midrule
\multirow{4}{*}{\textbf{Notifications}}
  & Notifications push FCM                           & \cmark \\
  & Centre de notifications in-app                   & \cmark \\
  & Badge sur l'onglet de navigation                 & \cmark \\
  & Configuration des abonnements (thème/auteur)     & \xmark \\
\midrule
\multirow{3}{*}{\textbf{Profil}}
  & Affichage des informations utilisateur           & \cmark \\
  & Galerie des photos publiées par l'utilisateur    & \xmark \\
  & Modification du profil                           & \xmark \\
\bottomrule
\multicolumn{3}{l}{\footnotesize \cmark~= Réalisé \quad \pmark~= Partiellement réalisé \quad \xmark~= Non encore réalisé}
\end{tabularx}
\end{table}

% ─────────────────────────────────────────────────────────────────────────────
\chapter{Points techniques remarquables}
% ─────────────────────────────────────────────────────────────────────────────

\section{Moteur de filtrage local (client-side)}

Pour éviter des allers-retours réseau à chaque frappe, tous les filtres sont appliqués
localement sur la liste complète des photos déjà chargée. La classe \texttt{FilterCriteria}
encapsule l'ensemble des critères actifs. La méthode \texttt{applyFilters()} du
\texttt{PhotoViewModel} itère sur cette liste et applique en cascade les critères de type,
d'auteur, de période (comparaison \texttt{Date}), de rayon (formule Haversine) et de texte
(correspondance dans lieu, pays, description, auteur, tags).

\section{Calcul de distance Haversine}

La formule de Haversine est utilisée pour calculer la distance en kilomètres entre deux
points GPS, ce qui permet le filtre \emph{autour d'un lieu} sans dépendance à un service
tiers~:

\begin{lstlisting}[language=Java, caption=Calcul de distance Haversine (PhotoViewModel)]
private double calculateDistance(double lat1, double lng1,
                                 double lat2, double lng2) {
    double R = 6371; // rayon moyen de la Terre en km
    double dLat = Math.toRadians(lat2 - lat1);
    double dLng = Math.toRadians(lng2 - lng1);
    double a = Math.sin(dLat/2)*Math.sin(dLat/2)
             + Math.cos(Math.toRadians(lat1))
             * Math.cos(Math.toRadians(lat2))
             * Math.sin(dLng/2)*Math.sin(dLng/2);
    double c = 2 * Math.atan2(Math.sqrt(a), Math.sqrt(1-a));
    return R * c;
}
\end{lstlisting}

\section{Configuration dynamique de l'URL du serveur}

L'écran \texttt{IpConfigActivity} permet de modifier à la volée l'adresse IP du backend
REST. Cette fonctionnalité est particulièrement utile lors des phases de test en réseau
local (e.g.\ machine virtuelle ou appareil physique sur le même sous-réseau Wi-Fi). La
nouvelle URL est passée à \texttt{RetrofitClient.setBaseUrl()}, ce qui force la
réinstanciation du client Retrofit.

\section{Adaptateur multi-vues avec bannière anonyme}

\texttt{PhotoAdapter} gère deux types d'items (\texttt{TYPE\_BANNER} et \texttt{TYPE\_PHOTO})
et trois modes d'affichage (grille, liste, card). La bannière est injectée en première
position pour inviter les utilisateurs non connectés à créer un compte. La
\texttt{GridLayoutManager.SpanSizeLookup} lui alloue automatiquement une largeur pleine
(span~= 2).

% ─────────────────────────────────────────────────────────────────────────────
\chapter{Fonctionnalités restantes et plan de travail}
% ─────────────────────────────────────────────────────────────────────────────

\section{Fonctionnalités prioritaires à compléter}

\begin{warningbox}[Travaux restants -- priorité haute]
\begin{enumerate}[leftmargin=*]
  \item \textbf{Annotations assistées par IA}~: enrichissement automatique de la photo
        lors de la publication (génération de tags, description courte) via un appel à
        une API d'IA (Claude, Gemini ou Vision API). C'est une fonctionnalité explicitement
        mentionnée dans la spécification.

  \item \textbf{Galerie de profil}~: affichage des photos publiées par l'utilisateur
        connecté dans son onglet Profil, avec possibilité de gestion (suppression).

  \item \textbf{Configuration des abonnements de notification}~: permettre à l'utilisateur
        de choisir les événements qui le notifient (publication d'un auteur donné, d'un
        groupe, d'un lieu ou d'un thème).

  \item \textbf{Signalement côté API}~: implémenter l'endpoint de signalement et connecter
        le bouton Signaler de \texttt{PhotoDetailActivity} à cet endpoint.

  \item \textbf{Filtre par similarité}~: la logique de récupération des photos similaires
        à une photo de référence (champ \texttt{similarPhotoId} déjà présent dans
        \texttt{FilterCriteria}) doit être branchée côté ViewModel et API.
\end{enumerate}
\end{warningbox}

\section{Améliorations envisagées}

\begin{itemize}
  \item \textbf{Pagination incrémentale}~: remplacer le chargement en une seule page
        (limite fixe à 100) par un défilement infini (\texttt{RecyclerView} +
        \texttt{EndlessScrollListener}).
  \item \textbf{Modification du profil}~: formulaire d'édition du nom, du pseudo et de
        l'avatar.
  \item \textbf{Mode hors-ligne partiel}~: mise en cache des dernières photos consultées
        (Room ou fichiers JSON dans le stockage interne).
  \item \textbf{Intégration TravelPath (Partie~3)}~: ajout du module passerelle permettant
        de naviguer entre TravelShare et TravelPath dans une navigation unifiée.
\end{itemize}

\section{Estimation de l'avancement global}

\begin{figure}[H]
\centering
\begin{tikzpicture}
  % Barre de progression
  \def\total{15}
  \def\done{11}
  \def\partial{2}
  \def\barW{12}

  % Fond
  \fill[lightgray, rounded corners=4pt] (0,0) rectangle (\barW, 0.8);
  % Réalisé
  \fill[success!70, rounded corners=4pt] (0,0) rectangle (\barW*\done/\total, 0.8);
  % Partiel
  \fill[warning!80] (\barW*\done/\total, 0) rectangle (\barW*(\done+\partial)/\total, 0.8);
  % Contour
  \draw[gray, rounded corners=4pt] (0,0) rectangle (\barW, 0.8);
  % Texte
  \node[anchor=west, font=\small\bfseries] at (0.2, 0.4) {\textcolor{white}{Réalisé}};
  \node[font=\small\bfseries] at (\barW*(\done+0.5*\partial)/\total, 0.4) {\textcolor{white}{\pmark}};
  \node[right, font=\small, color=darkgray] at (\barW + 0.2, 0.4)
        {\textbf{\done/\total} fonctionnalités (\pgfmathprintnumber[fixed]{\done*100/\total}\%)};
\end{tikzpicture}
\caption{Estimation visuelle de l'avancement (fonctionnalités principales)}
\label{fig:progress}
\end{figure}

Sur les \textbf{15~fonctionnalités principales} identifiées dans la spécification~:
\begin{itemize}
  \item \textbf{11} sont entièrement réalisées (\cmark),
  \item \textbf{2} sont partiellement réalisées (\pmark) : filtre par similarité et signalement,
  \item \textbf{2} restent à réaliser (\xmark) : annotations IA et configuration des abonnements de notification.
\end{itemize}

% ─────────────────────────────────────────────────────────────────────────────
\chapter{Conclusion}
% ─────────────────────────────────────────────────────────────────────────────

À mi-parcours, la partie \textbf{TravelShare} du projet \emph{Traveling} présente une base
fonctionnelle solide. L'architecture MVVM mise en place garantit une séparation claire des
responsabilités et une bonne maintenabilité. La couche réseau est opérationnelle et couvre
l'ensemble des endpoints nécessaires à ce stade. Les fonctionnalités cœur de l'application
-- parcours du fil de photos, recherche multicritère, publication, groupes et notifications
push -- sont toutes opérationnelles.

Les prochaines semaines seront consacrées principalement à l'intégration des annotations
assistées par IA, à l'enrichissement du profil utilisateur, à la finalisation des
fonctionnalités partielles (signalement, similarité) et à la préparation de la passerelle
d'intégration avec TravelPath.

\begin{successbox}[Points forts à mi-parcours]
\begin{itemize}[noitemsep]
  \item Architecture MVVM propre et extensible.
  \item Recherche vocale native (RecognizerIntent) fonctionnelle.
  \item Moteur de filtrage multi-critères entièrement côté client (réactif, sans latence réseau).
  \item Upload d'images vers Cloudinary intégré.
  \item Notifications push Firebase Cloud Messaging opérationnelles.
  \item Vue carte Google Maps avec marqueurs cliquables.
  \item Itinéraire depuis la fiche photo vers Google Maps.
\end{itemize}
\end{successbox}

\end{document}
